\begin{table}[H]
    \begin{minipage}[t]{.48\linewidth}
      \centering
      \caption*{ÁRBOL DE PROBLEMAS}
      
      \addstackgap[5pt]{\begin{tabularx}{\linewidth}{|l|X|}
        \hline
        \textbf{ID} & \textbf{Consecuencias} \\ \hline
        1 & Aumento en el tiempo de implementación \\ \hline
        2 & Mayor carga de trabajo para los profesionales SySO \\ \hline
        3 & Costos elevados para las empresas \\ \hline
        4 & Menos empresas cumplen con la normativa en materia de SySO \\ \hline
        5 & Baja satisfacción del profesional SySO \\ \hline
      \end{tabularx}}
    
    \justifying{
      \noindent \textbf{Problema Central}: 
        \begin{singlespace} 
            Los profesionales SySO de categoría A dedican mediante el método manual un 60\% del tiempo de desarrollo de un PGSST en llenado de formularios y planificación.
        \end{singlespace}
    }%justifying
    
      \noindent\addstackgap[5pt]{\begin{tabularx}{\linewidth}{|l|X|}
        \hline
        \textbf{ID} & \textbf{Causas Raíces del Problema} \\ \hline
        1 & Habilidad del profesional para transcribir datos \\ \hline
        2 & Instrumentos de medición que no permiten descargar sus datos en el formato requerido \\ \hline
        3 & Capacitación insuficiente del personal en el llenado de registros y plantillas \\ \hline
        4 & Adecuación de plantillas a casos particulares \\ \hline
      \end{tabularx}}
      
    \end{minipage}%
    \hfill % add space between minipages
    \begin{minipage}[t]{.48\linewidth}
      \centering
      \caption*{ÁRBOL DE OBJETIVOS}
      
      \addstackgap[5pt]{\begin{tabularx}{\linewidth}{|l|X|}
        \hline
        \textbf{ID} & \textbf{Resultados Esperados} \\ \hline
        1 & Reducción del tiempo de implementación de los PGSST \\ \hline
        2 & Menor carga de trabajo administrativo para el profesional SySO \\ \hline
        3 & Reducción de costos en elaboración de PGSST \\ \hline
        4 & Aumento en el cumplimiento de la normativa en materia de SySO \\ \hline
        5 & Aumento en la satisfacción del profesional SySO \\ \hline
      \end{tabularx}}
        
    \justifying{
      \noindent \textbf{Objetivo General:} 
        \begin{singlespace}
            Desarrollar un prototipo de software para disminuir el tiempo que los profesionales SySO de categoría A invierten llenado de formularios y planificación  al elaborar un PGSST.
        \end{singlespace}
    }%justifying

      \noindent\addstackgap[5pt]{\begin{tabularx}{\linewidth}{|l|X|}
        \hline
        \textbf{ID} & \textbf{Objetivos Específicos} \\ \hline
        1 & Diseñar un interfaz gráfica reduciendo la dependencia de habilidades específicas.\\ \hline
        2 & Implementar una herramienta para la transcripción automática de datos desde los instrumentos de medición. \\ \hline
        %3 & Facilitar un chat bot en la aplicación para que los clientes del profesional SySO consulten respecto al llenado de plantillas. \\ \hline
        3 & Implementar un modulo de generación plantillas que se adapten al caso. \\ \hline
      \end{tabularx}}
        
    \end{minipage} 
    
\end{table}
